\begin{abstract}
山区洪涝灾害可能同时削弱道路通行和通信覆盖能力。无人机应急投送因而需要统筹物资组批、地形净空、能源周转、交付时限与连续通信，单独缩短航程并不足以保证任务可执行。围绕给定标准场景，本文建立从单点运输能力到运输通信协同、再到任务分区的递进建模框架。

对于单点直接往返任务，根据载荷相关航程和爬升能耗定义最大安全载荷，采用可行子集枚举与集合划分动态规划描述不可拆货箱组批，并明确架次数、能耗和累计作业时间的优先关系。对于多点多架次运输，将逐站交接时刻、实体无人机占用和共享电池充电纳入统一的事件递推模型，以硬时限可行为前提，比较配送及时性与任务完成时间。

对于通信约束，依据双向链路预算和地形遮挡建立直连及单中继服务判据，将悬停点、服务区间、能源组件和运输时序联合考虑；连续通信以全时段可用为约束，离散采样仅作为辅助检查。对于任务分区，以同架次共访关系构造不可拆任务块，再依据实际通信分段确定各组的中继需求，在保持原任务时序的条件下通过资源占用区间计算独立配置规模。

附件统计表明，标准场景包含15个服务区、80个货箱，物资总质量为758 kg，总体积为2.011 m$^3$；医疗及首批保障要求合并后形成31个硬时限货箱，其余49箱采用软迟到评价。本文据此给出数据特征、公共物理模型、四问求解方法及验证设计。现阶段尚未形成经过完整复验的调度与分区定量结论，最大安全载荷、最终路线、通信余量及资源缺口须在统一物理与保障关系复核后确定。

\keywords{应急物资运输；不可拆货箱组批；共享电池调度；连续通信；任务分区}
\end{abstract}
